\section{Finite constants}\label{sec:constants}\label{app:coefficient-finite}
The main text writes most constants as $C(K)$ or $O_K(\cdot)$. Here we
make them explicit: each is an absolute number times a power of $K+1$
and a factor $e^{bK}$, which gives explicit degree thresholds and lets
Section~\ref{sec:log-rate} take the width $K$ of order $\log N$.
Let $C_*$ be the positive universal constant in
Nazarov--Nishry--Sodin, Corollary~1.2, and let $c_{\rm HW}>0$ be the
constant in Rudelson--Vershynin, Theorem~1.1.
We keep the dependence on both constants in all the bounds below.
We choose, for fixed parameters, the least index beyond which every
displayed inequality holds, so that all degree cutoffs are deterministic.
Table~\ref{tab:const-summary} lists the degree thresholds and the
exceptional probability of each result of the main text, and
Table~\ref{tab:named-constants} lists the named constants.

\begin{table}[!ht]
\centering
\small
\setlength{\tabcolsep}{4pt}
\begin{tabular}{@{}>{\raggedright\arraybackslash}p{0.2\textwidth}
>{\raggedright\arraybackslash}p{0.46\textwidth}
>{\raggedright\arraybackslash}p{0.3\textwidth}@{}}
\toprule
Result & Degree thresholds & Exceptional probability\\
\midrule
\ref{app:const-profile}, Lemma~\ref{lem:profile-rate}
 & $N\ge2K$ & none\\
\ref{app:const-log}, Theorem~\ref{thm:sign-logarithms}
 & $K\ge2$ and $N\ge(160e^{8K})^2$, uniformly for $r\in[1-K/N,1+K/N]$
 & $p_K^*(N)$\\
\ref{app:const-radial}, Corollary~\ref{cor:radial-probability}
 & those of \ref{app:const-log} and $N\ge2K$
 & $4p_K^*(N)$ for \eqref{eq:layer2-radial}, $5p_K^*(N)$ for
 \eqref{eq:radial-count-near-unit}, $5p_2^*(N)$ for its form at $K=2$\\
Lemma~\ref{lem:local-counts}
 & $t_N\le1$ and $N\ge e$ & $N^{-3}$\\
Lemma~\ref{lem:second-derivative}
 & $N\ge\max\{2,(K+1)/3\}$ & $N^{-10}$\\
\ref{app:const-small-derivatives}, Lemma~\ref{lem:small-derivatives}
 & $N\ge(6400e^{8K})^2$, and those of Lemma~\ref{lem:local-counts}
 and \eqref{eq:localization-thresholds}
 & the right-hand side of \eqref{eq:layer2-bad}\\
Lemma~\ref{lem:tails}
 & $j\ge60$ and $N\ge\max\{4,K\}$ & $N^{-10}$\\
Lemma~\ref{lem:isolation-interface}
 & $4aL/N^{3/2}\le1/N$ and $4S_KaL^2N^{-1/2}\sqrt{\log N}\le1$,
 that is, \eqref{eq:localization-thresholds} & none\\
\ref{app:block-bounds}, Section~\ref{sec:layer3}
 & $R\ge|x|+2$, $K\ge R$, $N\ge2R$ and $N\ge(160e^{8R})^2$, which implies
 $N>R$, those of the second to eighth rows at the
 width $K$, where Lemma~\ref{lem:tails} enters with $K+2$ in place of
 $K$ (so that $N\ge\max\{4,K+2\}$), and $j\ge60$
 & \eqref{eq:block-failure}, with $R=2$ and $x=0$ for
 Corollary~\ref{thm:stronglaw}\\
\ref{app:growing-annulus}, Section~\ref{sec:log-rate}
 & the conditions of Table~\ref{tab:growing-annulus}, $j\ge60$, $K\ge2$,
 $A\log N/N\le1$ and $N\ge K+2$
 & \eqref{eq:block-failure}, with $K=K_N$, $R=2$ and $x=0$\\
\ref{app:const-nearby}, Corollary~\ref{cor:crossings}
 & $0<\kappa_*<\min(\kappa,1/64)$ and
 $N^{-1-\kappa}+2a/d_0<N^{-1-\kappa_*}$
 & \eqref{eq:block-failure}, with $R=2$ and the thin-band events of
 \ref{app:const-radial}\\
\bottomrule
\end{tabular}
\caption{The degree thresholds and exceptional probabilities of the
results of the main text, with the subsection of
Appendix~\ref{sec:constants}, where there is one, that treats each. For other coefficient
distributions, Section~\ref{sec:model-constants} enlarges the constants
and adds exceptional probabilities.}
\label{tab:const-summary}
\end{table}

In Table~\ref{tab:named-constants}, $A$, $C_1(K)$, $S_K$ and $\lambda_K$
are the constants of Lemmas~\ref{lem:local-counts}
and~\ref{lem:second-derivative} and of \eqref{eq:jet-eigenvalue}, where
$S_K$ is written $C(K)$ and $\lambda_K$ is written $\lambda$.
The law constants $A_B$, $C_B$, $C_{1,G}$ and $C_{1,B}$ are in
Tables~\ref{tab:law-steinhaus}--\ref{tab:law-bounded}, and $A(B,\eta)$
and $C_1(K,B,\eta)$ in Lemma~\ref{lem:local-counts-general}.
The symbols $A$, $C_0$ and $D_0$ in Corollary~\ref{cor:T3-powers}, $B_0$
in Section~\ref{sec:layer2-models}, and $E_1,E_2$ in
Appendix~\ref{sec:general-estimates} have other meanings.
\begin{table}[p]
\centering
\footnotesize
\setlength{\tabcolsep}{3pt}
\caption{The named constants. Each is at most $C_0(K+1)^pe^{bK}$ with the
listed $(C_0,p,b)$, and the absolute constants have $p=b=0$.}
\label{tab:named-constants}
\begin{tabular}{@{}>{\raggedright\arraybackslash}p{0.1\textwidth}
>{\raggedright\arraybackslash}p{0.27\textwidth}
>{\raggedright\arraybackslash}p{0.13\textwidth}
>{\raggedright\arraybackslash}p{0.14\textwidth}
>{\raggedright\arraybackslash}p{0.28\textwidth}@{}}
\toprule
Constant & Value & $(C_0,p,b)$ & Defined & Used\\
\midrule
$D_K$ & $10(K+1)^2e^{4K}$, and $D_2=90e^8$ & $(10,2,4)$
 & \ref{app:const-profile}
 & \ref{app:const-radial}, \ref{app:block-bounds},
 \ref{app:growing-annulus}, \ref{sec:log-rate},
 Table~\ref{tab:rate-width}\\
$H_K$ & $10^8e^{20(K+1)}$, and $H_2=10^8e^{60}$ & $(H_0,0,20)$
 & \ref{app:const-log}
 & \ref{app:const-radial}, \ref{sec:model-constants},
 \ref{app:growing-annulus}, Table~\ref{tab:rate-width}\\
$C_1(K)$ & $(4A+4K+4)/\log2$, and $C_1(K+2)=(4A+4K+12)/\log2$
 & $(C_{10},1,0)$ for $C_1(K+2)$
 & Lemma~\ref{lem:local-counts}
 & \ref{sec:derivative-detection}, \ref{sec:layer2-models},
 \ref{sec:log-rate}, \ref{app:const-small-derivatives},
 \ref{sec:model-constants}, \ref{app:growing-annulus}\\
$S_K$ & $100e^{K+4}$ & $(S_0,0,1)$
 & Lemma~\ref{lem:second-derivative}
 & \ref{app:const-small-derivatives}, \ref{app:block-bounds},
 \ref{sec:model-constants}, \ref{app:growing-annulus}\\
$\lambda_K$ & $e^{-4K}/80$ & $(1/80,0,-4)$
 & \eqref{eq:jet-eigenvalue}
 & \ref{app:const-small-derivatives}, \ref{sec:model-constants},
 \ref{app:growing-annulus}\\
$D'_K$ & $2^{18}(K+1)S_K^2$ & $(D_0,1,2)$
 & \ref{app:const-small-derivatives}
 & \ref{sec:model-constants}, \ref{app:growing-annulus},
 Proposition~\ref{prop:sparse-exponents}\\
$\mathcal B_K$ & $2000e^{3K}(2/\lambda_K)^{3/2}$ & $(B_0,0,9)$
 & \ref{app:const-small-derivatives}
 & \ref{sec:model-constants}, \ref{app:growing-annulus},
 Proposition~\ref{prop:sparse-exponents}\\
$\mathcal T_K$ & $80e^{4K}/\lambda_K=6400e^{8K}$ & $(6400,0,8)$
 & \ref{app:const-small-derivatives}
 & \ref{sec:model-constants}, \ref{app:growing-annulus},
 Proposition~\ref{prop:sparse-exponents}\\
$V_K$ & $50+3\mathcal B_K+\mathcal T_K$ & $(V_0,0,9)$
 & \ref{app:const-small-derivatives}
 & \ref{sec:model-constants}, \ref{app:growing-annulus}\\
$c_K$ & $9/(1024S_K^2\lambda_K^2)$ & $(Q_0,0,6)$
 & \ref{app:const-small-derivatives}
 & \ref{sec:model-constants}, \ref{app:growing-annulus}\\
$C_K$ & $V_K(D'_K)^2$ & $(C_{\rm bad},2,13)$
 & \ref{app:const-small-derivatives}
 & \eqref{eq:layer2-bad}, \eqref{eq:block-failure},
 \ref{app:growing-annulus}, \ref{sec:log-rate},
 Table~\ref{tab:rate-width}\\
$Z_K$ & $C_1(K+2)(1+D'_Kc_K)$ & $(Z_0,2,8)$
 & \ref{app:const-small-derivatives}
 & \eqref{eq:layer2-bad}, \ref{app:block-bounds},
 \ref{app:growing-annulus}, \ref{sec:log-rate},
 Table~\ref{tab:rate-width}\\
$Z'_K$ & $C_1(K+2)(40+D'_K\mathcal B_K)$ & $(Z'_0,2,11)$
 & \ref{app:const-small-derivatives}
 & \eqref{eq:layer2-bad}, \ref{app:block-bounds},
 \ref{app:growing-annulus}, \ref{sec:log-rate},
 Table~\ref{tab:rate-width}\\
\midrule
$C_*$ & constant of \eqref{eq:layer2-nns}, $C_*\ge1$ &
 & \eqref{eq:layer2-nns}
 & \ref{app:const-log}, $E_1$, $E_2$, \ref{sec:model-constants}\\
$c_{\rm HW}$ & constant of \eqref{eq:HW-signs} &
 & \eqref{eq:HW-signs} & $A$, \ref{sec:model-constants}\\
$A$ & $\max(1,1024\pi/c_{\rm HW})$ &
 & Lemma~\ref{lem:local-counts}
 & $C_1(K)$, $C_{10}$, \ref{sec:log-rate}, \ref{app:growing-annulus}\\
$E_1$ & $\frac52+2e^{1/2}(4eC_*)^6$ &
 & \eqref{eq:layer2-log-split} & $E_0$\\
$E_2$ & $\frac1{16}+e^{1/2}(4eC_*)^6$ &
 & \eqref{eq:layer2-log-split} & $E_0$\\
$H_0$ & $10^8e^{20}$ & & \ref{app:growing-annulus}
 & $E_0$, \eqref{eq:E2-majorants}, \ref{sec:log-rate}\\
$S_0$ & $100e^4$ & & \ref{app:growing-annulus}
 & $D_0$, $Q_0$, Table~\ref{tab:growing-annulus}, \ref{sec:log-rate}\\
$D_0$ & $2^{18}S_0^2$ & & \ref{app:growing-annulus}
 & $C_{\rm bad}$, $Z_0$, $Z'_0$\\
$B_0$ & $2000\cdot160^{3/2}$ & & \ref{app:growing-annulus}
 & $V_0$, $Z'_0$\\
$V_0$ & $50+3B_0+6400$ & & \ref{app:growing-annulus} & $C_{\rm bad}$\\
$Q_0$ & $9\cdot80^2/(1024S_0^2)$ & & \ref{app:growing-annulus} & $Z_0$\\
$C_{10}$ & $(4A+12)/\log2$ & & \ref{app:growing-annulus}
 & $Z_0$, $Z'_0$\\
$C_{\rm bad}$ & $V_0D_0^2$ & & \ref{app:growing-annulus}
 & \eqref{eq:E2-majorants}, \ref{sec:log-rate}\\
$Z_0$ & $C_{10}(1+D_0Q_0)$ & & \ref{app:growing-annulus}
 & \eqref{eq:E2-majorants}, \ref{sec:log-rate}\\
$Z'_0$ & $C_{10}(40+D_0B_0)$ & & \ref{app:growing-annulus}
 & \eqref{eq:E2-majorants}, \ref{sec:log-rate}\\
$E_0$ & $E_1+E_2H_0$ & & \ref{app:growing-annulus}
 & \eqref{eq:E2-majorants}, \ref{sec:log-rate}\\
\bottomrule
\end{tabular}
\end{table}

\subsection{The variance profile}\label{app:const-profile}
For the constant of Lemma~\ref{lem:profile-rate} we set, for an integer
$K\ge2$,
\[
 D_K=10(K+1)^2e^{4K}.
\]
For $N\ge2K$ and $|x|\le K$, Lemma~\ref{lem:profile-rate} and
$1\le9(K+1)^2$ give
\begin{align*}
 \left|\log\sigma_N(1+x/N)-\tfrac12\log N-F(x)\right|
 &\le\frac{K^2+e^{4K}/2}{N}
 \le\frac{(K+1)^2e^{4K}+e^{4K}}{N}\\
 &\le\frac{10(K+1)^2e^{4K}}N=\frac{D_K}N.
\end{align*}

\subsection{Concentration of the angular logarithm}\label{app:const-log}
The one- and two-point bounds give the occupation constant $H_K$ of
Theorem~\ref{thm:sign-logarithms}, and Lemma~\ref{lem:clipping} gives
its error $\varepsilon_K^*(N)$ and exceptional probability $p_K^*(N)$.
Let $K\ge2$ be an integer, set
\[
 H_K=10^8e^{20(K+1)},
\]
and recall that $T=(\log N)/32$ is the clipping level of
Lemma~\ref{lem:clipping}.
The proof of Theorem~\ref{thm:sign-logarithms} uses $N\ge2K$, in
\eqref{eq:layer2-abel} and \eqref{eq:radial-bounds}, $N\ge(160e^{8K})^2$,
so that $40e^{8K}N^{-1/2}\le1/4$ for the covariance matrices of
Section~\ref{sec:h1-rademacher}, and $N\ge e$, in
Lemma~\ref{lem:clipping}. All three hold for $N\ge(160e^{8K})^2$,
since $(160e^{8K})^2=25600e^{16K}$ exceeds $2K$ and $e$.

Let $K\ge2$, $N\ge(160e^{8K})^2$ and $r\in[1-K/N,1+K/N]$.
In \eqref{eq:point-bounds}, the line before the last is $N^{-1/2}$ times
$(42d^{1/4}+16)\cdot46e^{9K}+320e^{8K}$, where the dimension is
$d\in\{2,4\}$, so that $d^{1/4}\le4^{1/4}$. Since
$42\cdot4^{1/4}+16=42\sqrt2+16<75.4$,
$e^{8K}=e^{-K}e^{9K}\le e^{-2}e^{9K}$, $75.4\cdot46=3468.4$ and
$320e^{-2}<44$, we have
\begin{align*}
 (42\cdot4^{1/4}+16)\cdot46e^{9K}+320e^{8K}
 &<75.4\cdot46e^{9K}+320e^{-2}e^{9K}\\
 &<3469e^{9K}+44e^{9K}<4000e^{9K}.
\end{align*}
Hence the one- and two-point differences in \eqref{eq:point-bounds} are
at most $4000e^{9K}N^{-1/2}$ on the retained angles and pairs. As in
\eqref{eq:occupation-bounds}, where the removed sets have measures at
most $2N^{-1/2}$ and $10N^{-1/2}$, we get, uniformly for $t\ge0$,
\[
 |\mathbb EA_t-F_G(t)|\le(2+4000e^{9K})N^{-1/2},\qquad
 |\mathbb EA_t^2-F_G(t)^2|\le(10+4000e^{9K})N^{-1/2}.
\]
Since $12000e^{9K}+14\le12014e^{9K}<10^8e^{20(K+1)}=H_K$ and
$F_G(t)=1-e^{-t^2}$, these bounds, the variance identity of
\eqref{eq:occupation-bounds} and $0\le\mathbb EA_t+F_G(t)\le2$ give
\begin{align}
 |\mathbb EA_t-(1-e^{-t^2})|&\le(2+4000e^{9K})N^{-1/2}
 \le H_KN^{-1/2},\nonumber\\
 \operatorname{Var}(A_t)
 &\le|\mathbb EA_t^2-F_G(t)^2|+2|\mathbb EA_t-F_G(t)|\nonumber\\
 &\le(10+4000e^{9K})N^{-1/2}+2(2+4000e^{9K})N^{-1/2}\nonumber\\
 &=(14+12000e^{9K})N^{-1/2}\le H_KN^{-1/2}.
 \label{eq:layer2-occupation}
\end{align}
We take the deviation for the truncated logarithmic integral, the
parameter $u$ of Theorem~\ref{thm:T1}, to be $N^{-1/16}(\log N)^7$, and
we set
\begin{align*}
 \varepsilon_K^*(N)&=N^{-1/16}(\log N)^7+2H_KTN^{-1/2}\\
 &\quad+e^{1/2}(4eC_*)^6(2N^{-1/16}+H_KN^{-1/2})(\log N)^7+\tfrac32N^{-1/16},\\
 p_K^*(N)&=H_KN^{-3/8}+N^{-12}
 +4H_KT^2N^{-3/8}(\log N)^{-14}.
\end{align*}
Then
\begin{equation}
 \begin{gathered}
 \Prob\{|J_N(r)-\log\sigma_N(r)+\gamma/2|>\varepsilon_K^*(N)\}
 \le p_K^*(N),\\
 \varepsilon_K^*(N)=O_K(N^{-1/16}(\log N)^7).
 \end{gathered}
                                                        \label{eq:layer2-log}
\end{equation}
The first line of \eqref{eq:layer2-log} is Lemma~\ref{lem:clipping} with
$H=H_K$, $C=C_*$, $E=\Omega$ and $B_E=1$, which applies for
$N\ge(160e^{8K})^2$ by \eqref{eq:layer2-occupation}, \eqref{eq:layer2-nns}
and \eqref{eq:layer2-energy}. Its error is $\varepsilon_K^*(N)$, since
$B_E/2+1=\tfrac32$, and its probability bound is $p_K^*(N)$, since
$4T^2=(\log N)^2/256$.
The second line follows from the bound on the error in
Lemma~\ref{lem:clipping}, which at $C=C_*$ and $B_E=1$ reads
\begin{equation}
 \varepsilon_K^*(N)\le E_1N^{-1/16}(\log N)^7+E_2H_KN^{-1/2}(\log N)^7,
 \qquad
 \begin{aligned}
 E_1&=\tfrac52+2e^{1/2}(4eC_*)^6,\\
 E_2&=\tfrac1{16}+e^{1/2}(4eC_*)^6,
 \end{aligned}
 \label{eq:layer2-log-split}
\end{equation}
where $E_1$ and $E_2$ are absolute constants, since $C_*$ is. Only the
second term depends on $K$, and it carries the factor $N^{-1/2}$.

\subsection{Annular mass and thin bands}\label{app:const-radial}
Next we compute the constants of
Corollary~\ref{cor:radial-probability}.
As in its proof, the four secant radii $1\pm K/N$ and $1\pm K/(2N)$ lie
in $[1-K/N,1+K/N]$. Hence
their variance errors are at most $D_K/N$ for $N\ge2K$, by
Section~\ref{app:const-profile}, and each logarithmic error is at most
$\varepsilon_K^*(N)$, outside probability $p_K^*(N)$, for
$N\ge(160e^{8K})^2$, by \eqref{eq:layer2-log}.
Lemma~\ref{lem:annular-mass} at $\eta=1/2$, with
$\varepsilon=\varepsilon_K^*(N)$ and $D=D_K$, then applies outside the
four logarithmic exceptional events, whose union has probability at most
$4p_K^*(N)$, and gives
\begin{equation}
 \frac{\#\{\alpha:||\alpha|-1|>K/N\}}N
 \le\frac{2\log2}{K}+\frac8K\bigl(\varepsilon_K^*(N)+D_K/N\bigr).
 \label{eq:layer2-radial}
\end{equation}

We use all five radii $1+qN^{-1-\kappa}$, $q=-2,\dots,2$, for
$0<\kappa\le1/64$, and $N\ge(160e^{8K})^2$ implies the condition $N\ge4$
of Lemma~\ref{lem:thin-band}. By that lemma with
$\varepsilon=\varepsilon_K^*(N)$, whose hypothesis holds outside the five
events of \eqref{eq:layer2-log} at these radii, for $q=-1,0,1$ the same
sample satisfies
\begin{equation}
 \left|\frac{\nu_N(1+qN^{-1-\kappa})}N-\frac12\right|
 \le2N^{-\kappa}+4\varepsilon_K^*(N)N^\kappa
 \label{eq:radial-count-near-unit}
\end{equation}
outside a union of probability at most $5p_K^*(N)$.
The probabilities of the exceptional events are summable along $N=j^8$,
since $4T^2(\log N)^{-14}=(\log N)^{-12}/256\le1$ for $N\ge e$, so that
the three terms of $p_K^*(j^8)$ are at most $H_Kj^{-3}$, $j^{-96}$ and
$H_Kj^{-3}$.

At $K=2$ we have $H_2=10^8e^{20(2+1)}=10^8e^{60}$, and
$\varepsilon_2^*(N)$ and $p_2^*(N)$ are $\varepsilon_K^*(N)$ and
$p_K^*(N)$ with $10^8e^{60}$ in place of $H_K$, with the degree
threshold $N\ge(160e^{16})^2$. Since the five radii lie in
$[1-2/N,1+2/N]$, \eqref{eq:radial-count-near-unit} holds at $K=2$, with
$\varepsilon_2^*(N)$, outside probability $5p_2^*(N)$.
This form is used in the proof of Corollary~\ref{cor:crossings}.

\subsection{Few small derivatives}\label{app:const-small-derivatives}
We now compute the constants of Lemma~\ref{lem:small-derivatives},
following its proof at the annular width $K$, with the local-count and
tail lemmas at the width $K+2$.

Let $S_K=100e^{K+4}$ be the constant of
Lemma~\ref{lem:second-derivative}, so that
\eqref{eq:sign-second-derivative} holds with $C(K)=S_K$, and let
\[
 C_1(K+2)=\frac{4A+4K+12}{\log2}
\]
be the constant $C_1(K)$ of Lemma~\ref{lem:local-counts} at the width
$K+2$. We set
\begin{align*}
 \lambda_K&=e^{-4K}/80,&
 D'_K&=2^{18}(K+1)S_K^2,\\
 \mathcal B_K&=2000e^{3K}(2/\lambda_K)^{3/2},&
 \mathcal T_K&=80e^{4K}/\lambda_K,\\
 V_K&=50+3\mathcal B_K+\mathcal T_K,&
 c_K&=\frac9{1024S_K^2\lambda_K^2},\\
 C_K&=V_K(D'_K)^2,&
 Z_K&=C_1(K+2)(1+D'_Kc_K),\\
 & & Z'_K&=C_1(K+2)(40+D'_K\mathcal B_K).
\end{align*}
Here $\mathcal B_K$ is the Rai\v{c} bound in the proof of
Lemma~\ref{lem:small-derivatives}, which bounds the normal-approximation
error and the one-point error, and $\mathcal T_K$ bounds the Gaussian
comparison error.

In that proof, the mesh constant must be at most $1/(2S_K)$ for the
bounds \eqref{eq:annular-mesh-test} and at most one for the mesh count. Both conditions hold
for the value $1/(64S_K)$, since $S_K\cdot1/(64S_K)=1/64\le1/2$ and
$S_K\ge1$. With $L=N^{1/64}$, let
\[
 h=\frac1{64S_KNL\sqrt{\log N}}.
\]
We use exactly $\lceil2K/(Nh)\rceil+1$ equally spaced radial rows,
including both endpoints, and $\lceil8\pi/h\rceil$ angles on each row.
This mesh has spacing $h$, and by the count of mesh points in the proof
of Lemma~\ref{lem:small-derivatives}, with $1/(64S_K)$ as its mesh
constant, and $53\cdot64^2=217088<262144=2^{18}$, its total size
satisfies
\begin{align*}
 m&\le53(K+1)(64S_K)^2NL^2\log N=53\cdot64^2(K+1)S_K^2NL^2\log N\\
 &\le2^{18}(K+1)S_K^2NL^2\log N=D'_KNL^2\log N.
\end{align*}
Since $Nh=1/(64S_KL\sqrt{\log N})$,
$N^2h^2\sqrt{\log N}=1/(4096S_K^2L^2\sqrt{\log N})$ and
$1/64+1/8192<1/16$, in \eqref{eq:annular-mesh-test}, with
$C(K)=S_K$ and the mesh constant $1/(64S_K)$, we have
\[
 \frac{Nh}L+\frac{S_K}2N^2h^2\sqrt{\log N}
 =\Bigl(\frac1{64}+\frac1{8192}\Bigr)\frac1{S_KL^2\sqrt{\log N}}
 \le\frac1{16S_KL^2\sqrt{\log N}},
\]
and $2(1+S_K/(64S_K))/L=2(1+1/64)/L=65/(32L)<3/L$.
It follows that the thresholds of Theorem~\ref{thm:T3} at $\tau=h$,
$R=2$ and $M_2=S_KN^{5/2}\sqrt{\log N}$ satisfy
\begin{equation}
 u\le(16S_K)^{-1}L^{-2}(\log N)^{-1/2},
 \qquad v_0\le3/L.
 \label{eq:layer2-meshtest}
\end{equation}
The proof of Lemma~\ref{lem:small-derivatives} uses $N\ge2K$ (in
\eqref{eq:layer2-abel} and for the radial weights), $N\ge K$ (for
$|z|\le2$), $N\ge4$ (for the retained angles on each row), $N\ge e$ (for
$h\le1/N$) and $N^{1/2}\ge6400e^{8K}$ (for the covariance lower bounds
$\lambda_K$ in dimension four and $\lambda_K/2$ in dimension eight).
Since $6400^2e^{16K}$ exceeds $2K$, $4$ and $e$, all of these hold if
\[
 N\ge(6400e^{8K})^2.
\]

By the proof of Lemma~\ref{lem:small-derivatives},
Theorem~\ref{thm:T3} applies with these thresholds, with $p_2=N^{-10}$,
$q_0=C_1(K+2)\log N$, $b_0=40C_1(K+2)\sqrt N\log N$ and $p_0=N^{-3}$, and
with $\lambda=\lambda_K$, $\delta=\delta_*=\mathcal B_KN^{-1/2}$,
$\chi=\mathcal T_KN^{-1/2}$ and $b_*=10N^{-1/2}$.
Since $u^2\le(16S_K)^{-2}L^{-4}(\log N)^{-1}$ and $v_0^2\le9L^{-2}$ by
\eqref{eq:layer2-meshtest}, we have
\[
 \frac{u^2v_0^2}{4\lambda_K^2}
 \le\frac1{4\lambda_K^2}\cdot\frac9{256S_K^2}\cdot\frac{L^{-6}}{\log N}
 =\frac{c_KL^{-6}}{\log N},
\]
so that the line
$u^2v_0^2/(4\lambda^2)+2000e^{3K}(2/\lambda)^{3/2}N^{-1/2}$ of
\eqref{eq:annular-jet-small-ball} gives at $\lambda=\lambda_K$ the
one-point probability
\begin{equation}
 c_KL^{-6}/\log N+\mathcal B_KN^{-1/2}.
 \label{eq:layer2-smallball}
\end{equation}
For a retained separated pair, the covariance of the two indicators in
the proof of Theorem~\ref{thm:T3} is at most
$\delta_*+\chi+2\delta=(3\mathcal B_K+\mathcal T_K)N^{-1/2}$, and
$V_*=(10+3\mathcal B_K+\mathcal T_K)N^{-1/2}\le V_KN^{-1/2}$. Hence, by
\eqref{eq:layer2-smallball} and the bound on $m$,
\begin{equation}
 \mu_*\le D'_K\{c_KNL^{-4}+\mathcal B_KN^{1/2}L^2\log N\},\qquad
 m^2V_*\le V_K(D'_K)^2N^{3/2}L^4(\log N)^2.
 \label{eq:layer2-Z}
\end{equation}
We take $t=NL^{-4}$ in \eqref{eq:T3-general}. By \eqref{eq:layer2-Z}
and $\sqrt N\le N^{1/2}L^2\log N$ for $N\ge e$, we keep the terms with
$NL^{-4}$ apart from those with $N^{1/2}L^2$:
\begin{align*}
 b_0+q_0(\mu_*+t)
 &\le C_1(K+2)(\log N)\bigl[40\sqrt N+NL^{-4}\\
 &\qquad+D'_K\{c_KNL^{-4}+\mathcal B_KN^{1/2}L^2\log N\}\bigr]\\
 &\le Z_KNL^{-4}\log N+Z'_KN^{1/2}L^2(\log N)^2.
\end{align*}
Since $t^2=N^2L^{-8}$ and $L^{12}=N^{3/16}$, \eqref{eq:layer2-Z} gives
\[
 \frac{m^2V_*}{t^2}
 \le V_K(D'_K)^2N^{-1/2}L^{12}(\log N)^2
 =C_KN^{-5/16}(\log N)^2,
\]
and hence, by \eqref{eq:T3-general}, with $p_0=N^{-3}$ the probability
of the exceptional local-count event of Lemma~\ref{lem:local-counts} and
$p_2=N^{-10}$ that of the derivative event of
Lemma~\ref{lem:second-derivative},
\begin{equation}
 \begin{aligned}
 &\Prob\{B_{K,N}>Z_KNL^{-4}\log N+Z'_KN^{1/2}L^2(\log N)^2\}\\
 &\qquad\le N^{-3}+N^{-10}+C_KN^{-5/16}(\log N)^2.
 \end{aligned}
 \label{eq:layer2-bad}
\end{equation}
In this union, the local and derivative events are counted once.
Since $NL^{-4}=N^{15/16}$, $N^{1/2}L^2=N^{17/32}$ and
$N^{17/32}\log N\le N^{15/16}$, the count $B_{K,N}$ is at most
$Z_KN^{15/16}\log N+Z'_KN^{17/32}(\log N)^2\le(Z_K+Z'_K)N^{15/16}\log N$
outside this union. The constant $Z_K$ of the leading term grows like
$(K+1)^2e^{8K}$, and $Z'_K$ like $(K+1)^2e^{11K}$
(Table~\ref{tab:named-constants}). No further threshold on $N$ is
needed here beyond $N\ge(6400e^{8K})^2$ above, the
local-count thresholds $t_N\le1$ and $N\ge e$ of
Lemma~\ref{lem:local-counts}, and the localization thresholds
\eqref{eq:localization-thresholds}.

\subsection{The radial law}\label{app:block-bounds}\label{app:const-moving-radii}
In this subsection we write out in finite form the bounds of the proof of
Theorem~\ref{thm:radial-profile} in Section~\ref{sec:layer3}, and hence
of Corollary~\ref{thm:stronglaw}, which is its case $x=0$ with $R=2$.
Recall that $N=j^8$, $m_j=(j+1)^8-j^8$, $L=N^{1/64}$, $h=N^{-1/64}$ and
$r_N=1+x/N$.
Fix $x$ and integers $R\ge|x|+2$ and $K\ge R$, and let $H_R$, $D_R$,
$\varepsilon_R^*(N)$ and $p_R^*(N)$ be $H_K$, $D_K$, $\varepsilon_K^*(N)$
and $p_K^*(N)$ at the width $R$, with the degree threshold
$N\ge(160e^{8R})^2$.
At $R=2$ these are $\varepsilon_2^*(N)$ and $p_2^*(N)$ of
Section~\ref{app:const-radial}.

Let $a=15e^{2(K+2)}\sqrt{m_j\log N}$ be the tail amplitude of Step~1 of
Section~\ref{sec:layer3}, that is, that of Lemma~\ref{lem:tails} at the
width $K+2$ with $m_{\max}=m_j$. With $M_2=S_KN^{5/2}\sqrt{\log N}$ by
\eqref{eq:sign-second-derivative} and $d_0=N^{3/2}/L$, the localization
conditions $4a/d_0\le1/N$ and $4M_2a\le d_0^2$ of
Table~\ref{tab:unit-conditions} read
\begin{equation}
 4aL/N^{3/2}\le1/N,\qquad4S_KaL^2N^{-1/2}\sqrt{\log N}\le1,
 \label{eq:localization-thresholds}
\end{equation}
since
\[
 \frac{4a}{d_0}=\frac{4aL}{N^{3/2}},\qquad
 \frac{4M_2a}{d_0^2}=\frac{4S_KN^{5/2}\sqrt{\log N}\,aL^2}{N^3}
 =4S_KaL^2N^{-1/2}\sqrt{\log N}.
\]
The second condition implies the first: it gives
$a\le N^{15/32}/(4S_K\sqrt{\log N})$, and then $d_0=N^{95/64}$ gives
$4a/d_0\le N^{-65/64}/(S_K\sqrt{\log N})<1/N$.

By the first line of the bound on $U_N$ in Step~3 of
Section~\ref{sec:layer3}, by \eqref{eq:layer2-radial}, multiplied by
$N$, and by \eqref{eq:layer2-bad}, on the radial and derivative events the
number $U_N$ of unselected roots satisfies
\begin{align}
 U_N&\le\#\{\alpha:f_N(\alpha)=0,\ ||\alpha|-1|>K/N\}+B_{K,N}\nonumber\\
 &\le Z_KN^{15/16}\log N+Z'_KN^{17/32}(\log N)^2\nonumber\\
 &\quad+N\left[\frac{2\log2}{K}+\frac8K\{\varepsilon_K^*(N)+D_K/N\}\right].
 \label{eq:block-unselected}
\end{align}

We apply the logarithmic estimates at $x+qh$, $q=-2,\dots,2$, and
Lemma~\ref{lem:local-radial-secant} at $x+qh$, $q=-1,0,1$, with $R$ in
place of $K$, logarithmic error $\eta_{\log}=\varepsilon_R^*(N)$, profile
error $\varepsilon=D_R/N$, $F'=\Phi$, $M_F=1/2$ and $B_F=1$, since
$0\le\Phi\le1$ and $0\le\Phi'=2\operatorname{Var}_x(t)\le1/2$ by Step~4
of Section~\ref{sec:layer3}, which also implies that
$|\Phi(x+qh)-\Phi(x)|\le|q|h/2$. For $N\ge2R$ we get
\[
 \left|\frac{\nu_N(r_N+qh/N)}N-\Phi(x)\right|
 \le\frac{|q|h}2+\frac h4+\frac RN
 +\frac{2(\varepsilon_R^*(N)+D_R/N)(1+R/N)}h.
\]
Since $|q|h/2+h/4\le2h$ for $|q|\le1$ and $2(1+R/N)\le4$, we obtain
\begin{equation}
 \begin{aligned}
 \left|\frac{\nu_N(r_N+qh/N)}N-\Phi(x)\right|
 &\le\delta_R^*(N),\qquad q=-1,0,1,\\
 \delta_R^*(N)&:=2h+\frac RN+\frac{4(\varepsilon_R^*(N)+D_R/N)}h.
 \end{aligned}\label{eq:E1-local}
\end{equation}
Here the exceptional probability is at most $5p_R^*(N)$, and
$\delta_R^*(N)\to0$. For Corollary~\ref{thm:stronglaw}, at $x=0$ and
$R=2$, this reads
\begin{equation}
 \begin{aligned}
 \left|\frac{\nu_N(1+qN^{-65/64})}N-\frac12\right|
 &\le\delta_2^*(N),\qquad q=-1,0,1,\\
 \delta_2^*(N)&=2N^{-1/64}+\frac2N
 +4\Bigl(\varepsilon_2^*(N)+\frac{D_2}N\Bigr)N^{1/64}.
 \end{aligned}
 \label{eq:block-center-crossings}
\end{equation}

By \eqref{eq:E1-local} with $q=\pm1$, as in \eqref{eq:moving-band}, the
strict band $||\alpha|-r_N|<h/N$ contains at most $2\delta_R^*(N)N$
roots.
This band contains the root of every selected disk whose closure meets
the closed annulus between the circles of radii $r_N$ and $1+x/n$, for
$N\le n\le N+m_j$, once $2a/d_0+|x|m_j/N^2<h/N$, as in Step~4 of
Section~\ref{sec:layer3}. At the width $K\ge R$ this follows from the
thresholds. Indeed, since $d_0=N^{3/2}/L=N^{95/64}$, and
$a\le N^{15/32}/(4S_K\sqrt{\log N})$ by the second condition of
\eqref{eq:localization-thresholds} with $L^2=N^{1/32}$, we have
\[
 N^{65/64}\cdot\frac{2a}{d_0}=2aN^{-15/32}
 \le\frac1{2S_K\sqrt{\log N}}<\frac12,
\]
since $S_K\sqrt{\log N}>1$, and hence, by $N/h=N^{65/64}$ and
\eqref{eq:unit-block-length},
\begin{align*}
 \frac Nh\Bigl(\frac{2a}{d_0}+\frac{|x|m_j}{N^2}\Bigr)
 &=N^{65/64}\cdot\frac{2a}{d_0}+|x|m_jN^{-63/64}\\
 &<\frac12+9|x|N^{-7/64}\le\frac12+\frac12=1,
\end{align*}
since
$18|x|<18R<3e^{7R/4}<160^{7/32}e^{7R/4}\le N^{7/64}$, by $|x|<R$,
$18R<3e^{7R/4}$ for $R\ge2$, $160^{7/32}>3$ and $N\ge(160e^{8R})^2$.
Therefore the number $c(r_N,1+x/n)$ of these disks is at most
$2\delta_R^*(N)N$, and the matching bound of Step~5 of
Section~\ref{sec:layer3} becomes
$|\nu_n(1+x/n)-\nu_N(r_N)|\le2\delta_R^*(N)N+U_N+m_j$.
The chain \eqref{eq:unit-block}, with $\delta_R^*(N)$ in place of
$\varepsilon_N$, then gives
\[
 \left|\frac{\nu_n(1+x/n)}n-\Phi(x)\right|
 \le3\delta_R^*(N)+\frac{U_N}N+\frac{2m_j}{N},
 \qquad N\le n\le N+m_j.
\]

Collecting the exceptional probabilities, we find that the common
exceptional event $E_{K,R,x,j}$ has probability at most
\begin{equation}
 4p_K^*(N)+5p_R^*(N)+N^{-3}+2N^{-10}
                 +C_KN^{-5/16}(\log N)^2.\label{eq:block-failure}
\end{equation}
This union consists of the four radial logarithmic events, used for
$U_N$, the five probe events at the width $R$, used in
\eqref{eq:E1-local}, and the local-count, derivative, mesh-deviation and
maximal-tail events. The shared local-count and derivative events are
counted once each, and a single maximal-tail event controls every degree
in the block. For Corollary~\ref{thm:stronglaw} we take $x=0$ and $R=2$.
For every fixed $K,R$, the series of \eqref{eq:block-failure} over
$N=j^8$ is finite, since its terms are, up to constants, the
probabilities of Table~\ref{tab:unit-events}, where $p_K^*(N)$ and
$p_R^*(N)$ are explicit forms of $\Pi_K(N)$ and $\Pi_R(N)$, and every
exponent of $j$ in the last column of that table is less than $-1$.

\subsection{Other coefficient distributions}\label{sec:model-constants}
This subsection gives the constants with which the other distributions
satisfy (H1)--(H5) in Proposition~\ref{prop:coefficient-transfer}, as
used in the proof of Theorem~\ref{thm:radial-profile} and
Corollary~\ref{thm:stronglaw} for these distributions at the end of
Section~\ref{sec:layer2-models}.
For each distribution we recompute the constants of the sign case from
the inputs that depend on it: its logarithmic-moment constant, its energy
bound, its local-count constant and, for unbounded coefficients, an
envelope on their size.
Tables~\ref{tab:law-steinhaus}, \ref{tab:law-gaussian}
and~\ref{tab:law-bounded} list these constants.

\begin{table}[!ht]
\centering
\small
\begin{tabular}{@{}ll@{}}
\toprule
Input & Value\\
\midrule
Logarithmic moments $(C,\beta)$ & $(C_*,6)$\\
Energy & $E=\Omega$, $B_E=1$, $p_E=0$\\
Occupation constant & $H_K$\\
Local count & $C_1(K+2)$, exceptional probability $N^{-3}$\\
All other constants and thresholds & those of the sign case\\
\bottomrule
\end{tabular}
\caption{Constants for Steinhaus coefficients.}
\label{tab:law-steinhaus}
\end{table}

\paragraph{Steinhaus coefficients.}
These coefficients have the constants of the sign case
(Table~\ref{tab:law-steinhaus}), since their amplitudes have modulus
one, so that conditional signs yield the same $C_*$ and $C_1(K+2)$.

\begin{table}[!ht]
\centering
\small
\begin{tabular}{@{}>{\raggedright\arraybackslash}p{0.27\textwidth}
>{\raggedright\arraybackslash}p{0.69\textwidth}@{}}
\toprule
Input & Value\\
\midrule
Logarithmic moments $(C,\beta)$ & $(16,6)$\\
Energy & $E=\{V\le B_E\}$, $B_E=2$,
 $p_E=\exp[-3N/(32e^{6K})]$\\
Occupation constant & $H_K$\\
Local count & $C_{1,G}(K+2)=C_1(K+2)+2/\log2
 =(4A+4(K+2)+6)/\log2=(4A+4K+14)/\log2$,
 exceptional probability $N^{-3}+\exp[-3(N+1)/32]$\\
Small derivatives & $Z_{K,G}=C_{1,G}(K+2)(1+D'_Kc_K)$,
 $Z'_{K,G}=C_{1,G}(K+2)(40+D'_K\mathcal B_K)$; $D'_K$, $c_K$,
 $\mathcal B_K$, $V_K$ and $C_K=V_K(D'_K)^2$ unchanged\\
Derivative and tails & amplitudes of the sign case,
 $\lceil N^4\rceil\le2N^4$ boundary samples, on the envelope event
 $\max_{k\le2N}|g_k|\le N$ of probability at least
 $1-2(2N+1)e^{-N^2/2}$\\
Thresholds & $N\ge(160e^{8K})^2$, $N\ge e^2$, $\rho\le2$, and those of
 the sign case\\
\bottomrule
\end{tabular}
\caption{Constants for standard real and complex Gaussian coefficients.}
\label{tab:law-gaussian}
\end{table}

\paragraph{Real and complex Gaussian coefficients.}
By the proof of Proposition~\ref{prop:coefficient-transfer}, with
Lemma~\ref{lem:local-counts-gauss} for the local count, standard complex
Gaussian coefficients have the constants of the real ones
(Table~\ref{tab:law-gaussian}).

The choice $(C,\beta)=(16,6)$ is valid, since $(16p)^p\le(16p)^{6p}$ for
$p\ge1$, so that, by \eqref{eq:layer2-gaussianlog}, the logarithmic
moments satisfy \eqref{eq:layer2-nns} with $16$ in place of $C_*$.
In the proof of Proposition~\ref{prop:coefficient-transfer}, the Gaussian
occupation constants $2+320e^{8K}$ and $14+960e^{8K}$ are at most $H_K$,
since $2+320e^{8K}<14+960e^{8K}\le974e^{9K}<10^8e^{20(K+1)}=H_K$, so that
\eqref{eq:layer2-occupation} holds for these coefficients.
Hence Lemma~\ref{lem:clipping} applies with $H=H_K$, $C=16$,
$E=\{V\le B_E\}$ and $B_E=2$, where $V=\int|X_r|^2\dd\mu$ as in
Proposition~\ref{prop:coefficient-transfer}, and $\Prob(E^c)\le p_E$ by
\eqref{eq:layer2-gaussianenergy}.

This is the computation of Section~\ref{app:const-log} with $C_*$
replaced by $16$, $B_E=1$ by $2$ and the term $p_E$ added: since
$4e\cdot16=64e$, $B_E/2+1=2$ and $4T^2=(\log N)^2/256$, it gives, for
$N\ge(160e^{8K})^2$ and $r\in[1-K/N,1+K/N]$,
\begin{align*}
 &\Prob\bigl\{|J_N(r)-\log\sigma_N(r)+\gamma/2|
 >N^{-1/16}(\log N)^7+2H_KTN^{-1/2}\\
 &\qquad\quad+e^{1/2}(64e)^6(2N^{-1/16}+H_KN^{-1/2})(\log N)^7+2N^{-1/16}\bigr\}\\
 &\quad\le H_KN^{-3/8}+N^{-12}
 +\frac{H_K}{256}N^{-3/8}(\log N)^{-12}+p_E.
\end{align*}
For these coefficients, the error and the probability on the two sides of
this display replace $\varepsilon_K^*(N)$ and $p_K^*(N)$ at each radius in
\eqref{eq:layer2-radial}, \eqref{eq:radial-count-near-unit},
\eqref{eq:block-unselected}, \eqref{eq:block-center-crossings},
\eqref{eq:block-failure} and \eqref{eq:E1-local}, and in the radial
thresholds of Section~\ref{app:const-radial}.

The bounds of Lemmas~\ref{lem:second-derivative} and~\ref{lem:tails} keep
the amplitudes of the sign case on the envelope event, by the proof of
Proposition~\ref{prop:coefficient-transfer}: with $\lceil N^4\rceil$
samples in place of $\lceil N^3\rceil$, the derivative bounds gain the
factor $N$ and the sample spacing loses it.
The thresholds $N\ge e^2$ and $\rho\le2$ imply the conditions $N\ge4$ and
$\rho\le2$ of the interpolation step in the proof of Lemma~\ref{lem:tails}.
In $Z_K$ and $Z'_K$ we replace $C_1(K+2)$ by $C_{1,G}(K+2)$, while the
other constants of Section~\ref{app:const-small-derivatives} remain upper
bounds, since the normal approximation is exact. On the common
coefficient-envelope event, the bound \eqref{eq:layer2-bad} is then
replaced by
\begin{multline*}
 \Prob\{B_{K,N}>Z_{K,G}NL^{-4}\log N+Z'_{K,G}N^{1/2}L^2(\log N)^2\}\\
 \le N^{-3}+\exp[-3(N+1)/32]+N^{-10}+C_KN^{-5/16}(\log N)^2.
\end{multline*}
All other exceptional probabilities are those of the sign case.

\begin{table}[!ht]
\centering
\small
\begin{tabular}{@{}>{\raggedright\arraybackslash}p{0.27\textwidth}
>{\raggedright\arraybackslash}p{0.69\textwidth}@{}}
\toprule
Input & Value\\
\midrule
Logarithmic moments $(C,\beta)$ & $(C_B,6)$, with
 $C_B=2\{C_*+\max(\tfrac12\log2,\log B)+1\}$\\
Energy & $E=\{V\ge1/2\}$, $B_E=B^2$,
 $p_H=\exp[-N/(2B^4e^{6K})]$\\
Occupation constant & $H_{K,B}=(1+B^3)H_K$\\
Derivative and mesh & $S_{K,B}=BS_K$, $D'_{K,B}=B^2D'_K$,
 $\mathcal B_{K,B}=B^3\mathcal B_K$, $c_{K,B}=c_K/B^2$; $\lambda_K$ and
 $\mathcal T_K$ unchanged\\
Tail amplitude & $Ba$, with $a=15e^{2(K+2)}\sqrt{m_j\log N}$\\
Local count & $A_B=\max(1,2048\pi B^2/c_{\rm HW})$,
 $C_{1,B}(K+2)=(4A_B+4(K+2)+4+\log B)/\log2$,
 exceptional probability $N^{-3}+\exp[-(N+1)/(2B^4)]$\\
Small derivatives & $V_{K,B}$, $C_{K,B}$, $Z_{K,B}$, $Z'_{K,B}$: the
 formulas for $V_K$, $C_K$, $Z_K$, $Z'_K$ with these constants and
 $C_{1,B}(K+2)$\\
Thresholds & $A_B\log N/N\le1$, $N\ge e^2$, and the enlarged thresholds
 below\\
\bottomrule
\end{tabular}
\caption{Constants for a bounded centrally symmetric distribution with
$\mathbb E|\xi_0|^2=1$ and $|\xi_0|\le B$.}
\label{tab:law-bounded}
\end{table}

\paragraph{Centrally symmetric bounded coefficients.}
Assume that the coefficient distribution is nondegenerate and centrally
symmetric, as defined in Section~\ref{par:central-symmetry}.
We normalize $\mathbb E|\xi_0|^2=1$ and assume $|\xi_0|\le B$, so that
$B\ge1$.
We use $V$, $E=\{V\ge1/2\}$, $J_N$ and $C_B$ as in
Proposition~\ref{prop:coefficient-transfer} and its proof, and the
constants of Table~\ref{tab:law-bounded}, where $a$ is the tail
amplitude of Step~1 of Section~\ref{sec:layer3}.
The local-count constant is that of Lemma~\ref{lem:local-counts-bounded}
with $K+2$ in place of $K$, and the conditional trace event adds
$\exp[-(N+1)/(2B^4)]$ to its exceptional probability.
Since $\mathcal B_{K,B}=B^3\mathcal B_K$, $2/\lambda_K=160e^{4K}$,
$(160e^{4K})^{3/2}=160^{3/2}e^{6K}$, $80e^{4K}/\lambda_K=6400e^{8K}$,
$D'_{K,B}=B^2D'_K$, $(D'_K)^2=2^{36}(K+1)^2S_K^4$ and
$S_K^4=10^8e^{4K+16}$, and since $D'_{K,B}c_{K,B}=D'_Kc_K$ and
$D'_{K,B}\mathcal B_{K,B}=B^5D'_K\mathcal B_K$,
\begin{align*}
 V_{K,B}&=50+3\mathcal B_{K,B}+\mathcal T_K
 =50+6000B^3e^{3K}(2/\lambda_K)^{3/2}+\frac{80e^{4K}}{\lambda_K}\\
 &=50+6000\cdot160^{3/2}B^3e^{9K}+6400e^{8K},\\
 C_{K,B}&=V_{K,B}(D'_{K,B})^2=B^4V_{K,B}(D'_K)^2\\
 &=2^{36}\cdot10^8B^4(K+1)^2e^{4K+16}V_{K,B},\\
 Z_{K,B}&=C_{1,B}(K+2)(1+D'_{K,B}c_{K,B})=C_{1,B}(K+2)(1+D'_Kc_K),\\
 Z'_{K,B}&=C_{1,B}(K+2)(40+D'_{K,B}\mathcal B_{K,B})
 =C_{1,B}(K+2)(40+B^5D'_K\mathcal B_K).
\end{align*}
For the occupation bounds, the constant $c_B$ of the proof of
Proposition~\ref{prop:coefficient-transfer} satisfies, as in
Section~\ref{app:const-log} with the factor $B^3$ on the first term only
and with $B\ge1$ in the last step,
\[
 c_B=(42\cdot4^{1/4}+16)\cdot46B^3e^{9K}+320e^{8K}
 <3469B^3e^{9K}+44e^{9K}<4000B^3e^{9K}.
\]
Hence, by $B^3e^{9K}\ge1$, $12014<10^8$ and $9K<20(K+1)$, the occupation
constants $2+c_B<14+3c_B$ of that proof satisfy
$14+3c_B<14+12000B^3e^{9K}\le12014B^3e^{9K}<B^3H_K<H_{K,B}$, so that
\eqref{eq:layer2-occupation} holds with $H_{K,B}$ in place of $H_K$.
Hence Lemma~\ref{lem:clipping} applies with $H=H_{K,B}$, $C=C_B$,
$E=\{V\ge1/2\}$ and $B_E=B^2$, by the moment bound on $E$ in
Proposition~\ref{prop:coefficient-transfer} and since $V\le B^2$.
With $B_E/2+1=B^2/2+1$ and $4T^2=(\log N)^2/256$, its error is
\begin{align*}
 \varepsilon_{K,B}^*(N)
 &=N^{-1/16}(\log N)^7+2H_{K,B}TN^{-1/2}\\
 &\quad +e^{1/2}(4eC_B)^6(2N^{-1/16}+H_{K,B}N^{-1/2})(\log N)^7
 +(B^2/2+1)N^{-1/16},
\end{align*}
and it gives
\begin{align*}
 &\Prob\{|J_N(r)-\log\sigma_N(r)+\gamma/2|>\varepsilon_{K,B}^*(N)\}
 \le H_{K,B}N^{-3/8}+N^{-12}\\
 &\qquad+\frac{H_{K,B}}{256}N^{-3/8}(\log N)^{-12}
 +\exp[-N/(2B^4e^{6K})],
\end{align*}
where the last term is the probability of $E^c$ from
\eqref{eq:layer2-boundedenergy}.
As in the real Gaussian case, $\varepsilon_{K,B}^*(N)$ and the right-hand
side of the last display replace $\varepsilon_K^*(N)$ and $p_K^*(N)$ at
each radius.
Similarly, by the proof of Proposition~\ref{prop:coefficient-transfer}
with the deviation $t=NL^{-4}$, the bound \eqref{eq:layer2-bad} is
replaced by
\begin{multline*}
 \Prob\{B_{K,N}>Z_{K,B}NL^{-4}\log N+Z'_{K,B}N^{1/2}L^2(\log N)^2\}\\
 \le N^{-3}+\exp[-(N+1)/(2B^4)]+N^{-10}+C_{K,B}N^{-5/16}(\log N)^2.
\end{multline*}
We require $A_B\log N/N\le1$ and $N\ge e^2$, and all previous thresholds
with the enlarged constants.
Since $S_{K,B}\cdot Ba=B^2S_Ka$, these are
\begin{align*}
 &N\ge(6400e^{8K})^2,\qquad N\ge2K,\\
 &4BaL/N^{3/2}\le1/N,\qquad
  4B^2S_KaL^2N^{-1/2}\sqrt{\log N}\le1,
\end{align*}
together with the thresholds $N\ge\max\{2,(K+1)/3\}$, $j\ge60$ and
$N\ge\max\{4,K+2\}$, which do not change.
These inequalities hold eventually for each fixed distribution, with
cutoffs that depend on $B$.

If $\Prob(\xi_0=0)>0$, the trailing-zero allowance $\lceil D\log N\rceil$,
with $D>4/|\log\Prob(\xi_0=0)|$, and the bound
$N^{-D|\log\Prob(\xi_0=0)|}$ on the probability of its exceptional event
at each degree are those of the proof of
Theorem~\ref{thm:radial-profile} and Corollary~\ref{thm:stronglaw} for
other distributions in Section~\ref{sec:layer2-models}, which also
treats the initial zero multiplicity and the actual degrees in
Corollary~\ref{cor:T4-two-radii}.

\subsection{Bounds for a growing annulus}\label{app:growing-annulus}
In this subsection we compute the constants of the proof in
Section~\ref{sec:log-rate}, at the width $K=K_N$, where
$K_N=\lfloor(\log N)/128-\log\log N\rfloor$ is the width of that
section, with $N=j^8$ and $L=N^{1/64}$ as before, and with the
local-count and tail lemmas at the width $K+2$.
The following bounds hold whenever the displayed deterministic conditions
hold.
Set
\begin{align*}
 H_0&=10^8e^{20},&S_0&=100e^4,&D_0&=2^{18}S_0^2,\\
 B_0&=2000\cdot160^{3/2},&V_0&=50+3B_0+6400,&
 Q_0&=\frac{9\cdot80^2}{1024S_0^2},\\
 A&=\max(1,1024\pi/c_{\rm HW}),&C_{10}&=(4A+12)/\log2,&
 C_{\rm bad}&=V_0D_0^2,\\
 Z_0&=C_{10}(1+D_0Q_0),&Z'_0&=C_{10}(40+D_0B_0),&
 E_0&=E_1+E_2H_0,
\end{align*}
where $E_1$ and $E_2$ are the absolute constants of
\eqref{eq:layer2-log-split}.
These constants are independent of $K,N$.
By the definitions we have, for $K\ge0$,
\begin{align*}
 H_K&=H_0e^{20K},&S_K&=S_0e^K,&D'_K&=D_0(K+1)e^{2K},\\
 \mathcal B_K&=B_0e^{9K},&\mathcal T_K&=6400e^{8K},&V_K&\le V_0e^{9K},\\
 c_K&=Q_0e^{6K},&C_1(K+2)&\le C_{10}(K+1),&
 C_K&\le C_{\rm bad}(K+1)^2e^{13K},\\
 Z_K&\le Z_0(K+1)^2e^{8K},&Z'_K&\le Z'_0(K+1)^2e^{11K}.
\end{align*}
Indeed, the first three hold since $H_K=10^8e^{20}e^{20K}$,
$S_K=100e^4e^K$ and $D'_K=2^{18}(K+1)S_K^2=2^{18}S_0^2(K+1)e^{2K}$. For the
others we use $1/\lambda_K=80e^{4K}$, so that $2/\lambda_K=160e^{4K}$ and
$\mathcal T_K=80e^{4K}\cdot80e^{4K}=6400e^{8K}$, and, for the bounds,
$e^{8K}\le e^{9K}$, $1\le(K+1)e^{bK}$, $4K\le(4A+12)K$ and
$C_1(K+2)=(4A+4(K+2)+4)/\log2$:
\begin{align*}
 \mathcal B_K&=2000e^{3K}(160e^{4K})^{3/2}=2000\cdot160^{3/2}e^{9K}
 =B_0e^{9K},\\
 V_K&=50+3B_0e^{9K}+6400e^{8K}\le(50+3B_0+6400)e^{9K}=V_0e^{9K},\\
 c_K&=\frac{9\cdot80^2e^{8K}}{1024S_0^2e^{2K}}=Q_0e^{6K},\\
 C_1(K+2)&=\frac{4A+4K+12}{\log2}
 \le\frac{(4A+12)(K+1)}{\log2}=C_{10}(K+1),\\
 C_K&=V_K(D'_K)^2
 \le V_0e^{9K}\cdot D_0^2(K+1)^2e^{4K}
 =C_{\rm bad}(K+1)^2e^{13K},\\
 Z_K&=C_1(K+2)(1+D'_Kc_K)
 \le C_{10}(K+1)\{1+D_0Q_0(K+1)e^{8K}\}\\
 &\le Z_0(K+1)^2e^{8K},\\
 Z'_K&=C_1(K+2)(40+D'_K\mathcal B_K)
 \le C_{10}(K+1)\{40+D_0B_0(K+1)e^{11K}\}\\
 &\le Z'_0(K+1)^2e^{11K}.
\end{align*}

These are the forms $C_0(K+1)^pe^{bK}$ of
Table~\ref{tab:named-constants}, and the following lemma evaluates such
a form at $K=K_N$.

\begin{lemma}[Constants at the growing width]\label{lem:growing-width}
Let $N\ge e$ and $K_N\ge0$, and let $X_K\le C_0(K+1)^pe^{bK}$ with
$C_0\ge0$, $p\ge0$ and $b\ge0$. Then
\[
 X_{K_N}\le C_0(K_N+1)^pN^{b/128}(\log N)^{-b},
\]
and if moreover $K_N+1\le\log N$, then
$X_{K_N}\le C_0(\log N)^pN^{b/128}(\log N)^{-b}$.
\end{lemma}

\begin{proof}
By the floor, $K_N\le(\log N)/128-\log\log N$, and hence, since $b\ge0$
and $\log N\ge1$,
\[
 e^{bK_N}\le N^{b/128}(\log N)^{-b}\le N^{b/128},\qquad b\ge0.
\]
This gives the first bound, and the second follows since
$0<K_N+1\le\log N$ and $p\ge0$.
\end{proof}

By Lemma~\ref{lem:growing-width} with the forms above, keeping the
factor $(K+1)^p$ and, only for $Z_K$, the factor $(\log N)^{-8}$, we get
for $N\ge e$
\begin{align}
 \varepsilon_K^*(N)&\le E_0N^{-1/16}(\log N)^7,\nonumber\\
 D_K/N&\le10(K+1)^2N^{-31/32},\nonumber\\
 p_K^*(N)&\le H_0N^{-7/32}+N^{-12}
 +\frac{H_0}{256}N^{-7/32}(\log N)^{-12},\nonumber\\
 C_KN^{-5/16}(\log N)^2
 &\le C_{\rm bad}(K+1)^2N^{-27/128}(\log N)^2,\nonumber\\
 Z_KN^{-1/16}\log N
 &\le Z_0(K+1)^2(\log N)^{-7},\nonumber\\
 Z'_KN^{-15/32}(\log N)^2
 &\le Z'_0(K+1)^2N^{-49/128}(\log N)^2.
                                                        \label{eq:E2-majorants}
\end{align}
By Lemma~\ref{lem:growing-width} with $b=20$, we have
$H_K=H_0e^{20K}\le H_0N^{5/32}$, and $5/32-1/2=-11/32$ and
$5/32-3/8=-7/32$. By \eqref{eq:layer2-log-split} and
$N^{-11/32}\le N^{-1/16}$, the first line follows from
\[
 \varepsilon_K^*(N)
 \le E_1N^{-1/16}(\log N)^7+E_2H_0N^{-11/32}(\log N)^7
 \le E_0N^{-1/16}(\log N)^7.
\]
The second line follows from
$D_K/N=10(K+1)^2e^{4K}N^{-1}\le10(K+1)^2N^{1/32-1}=10(K+1)^2N^{-31/32}$.
Since $4T^2=(\log N)^2/256$, the third line follows from
\begin{align*}
 p_K^*(N)&=H_KN^{-3/8}+N^{-12}
 +\frac{H_K}{256}N^{-3/8}(\log N)^{-12}\\
 &\le H_0N^{-7/32}+N^{-12}
 +\frac{H_0}{256}N^{-7/32}(\log N)^{-12}.
\end{align*}
By the bounds on $C_K$, $Z_K$ and $Z'_K$ above, and
Lemma~\ref{lem:growing-width} with $b=13$, $b=8$ and $b=11$, the last
three lines follow from $13/128-5/16=-27/128$,
$e^{8K}N^{-1/16}\le(\log N)^{-8}$ and $11/128-15/32=-49/128$.

The caps on $K/\log N$ in Table~\ref{tab:rate-width} are $\beta/b$ in
its first part and $(\beta-1/8)/b$ in its second, where $b$ is the
exponent of the named constant of the row in
Table~\ref{tab:named-constants}, plus $2$ for the factor
$e^{2(K+2)}=e^4e^{2K}$ of $a$ in the two rows with $a$, and $b=16$ for
$(6400e^{8K})^2$.

Each ratio in Table~\ref{tab:growing-annulus} tends to zero, and the
indicated condition holds once the ratio is at most one.
Indeed, at $K=K_N$ every ratio other than $2K/N$ is a constant times
powers of $K_N+1\le\log N$ and of $\log N$, times a negative power of
$N$, and $2K_N/N\le2(\log N)/N$. These rates are compared in
Section~\ref{sec:log-rate}.
\begin{table}[!htbp]
\centering
\small
\begin{tabular}{>{\raggedright\arraybackslash}p{0.32\textwidth}p{0.61\textwidth}}
\toprule
Condition & Sufficient ratio bound\\
\midrule
$(160e^{8K})^2\le N$, $(6400e^{8K})^2\le N$ & $6400^2N^{-7/8}$\\
$2K\le N$ & $2K/N$\\
Taylor localization inequality
 & $180S_0e^4(\log N)N^{-1/128}$\\
$2a/d_0<N^{-65/64}$
 & $90e^4\sqrt{\log N}\,N^{-1/64}$\\
\bottomrule
\end{tabular}
\caption{Sufficient ratio bounds for the degree conditions at the width
$K=K_N$.}
\label{tab:growing-annulus}
\end{table}
We require the last ratio to be strictly below one. The condition
$4a/d_0\le1/N$ follows from the Taylor localization inequality, as shown
after \eqref{eq:localization-thresholds}.

We obtain the ratios of the table by dividing each condition by its
right-hand side and substituting the bounds above. For the first row,
since $160<6400$ and by Lemma~\ref{lem:growing-width} with $b=16$,
\[
 \frac{(6400e^{8K})^2}N=6400^2e^{16K}N^{-1}\le6400^2N^{1/8-1}
 =6400^2N^{-7/8}.
\]
The last two rows, for the inequality
$4S_KaL^2N^{-1/2}\sqrt{\log N}\le1$ of Taylor localization in
\eqref{eq:localization-thresholds}, which is $4M_2a\le d_0^2$, and for
the strict band $N^{65/64}\cdot2a/d_0<1$, are computed in the proof of
Theorem~\ref{thm:log-rate} from \eqref{eq:unit-radius} and
\eqref{eq:unit-parameters}, with $180S_0e^4=18000e^8$ and the
exponents $7/16+1/32-1/2=-1/32$, $3/128-1/32=-1/128$,
$7/16-15/32=-1/32$ and $2/128-1/32=-1/64$. They use the bound
$m_j\le9N^{7/8}$ of \eqref{eq:unit-block-length} and the bound
$a\le45e^4e^{2K}N^{7/16}\sqrt{\log N}$ of \eqref{eq:unit-amplitude},
where $e^{2(K+2)}=e^4e^{2K}$.
The small-derivative lemma adds no ratio to the table, since its
thresholds are $N\ge(6400e^{8K})^2$, those of
Lemma~\ref{lem:local-counts} and \eqref{eq:localization-thresholds}.

We also impose
\[
 j\ge60,\qquad K\ge2,\qquad A\log N/N\le1,\qquad N\ge K+2.
\]
These conditions are eventual.
Beyond a fixed cutoff, the secant and probe radii belong to their
required annuli.
In this way we obtain all local-count, covariance, tail, localization
and radial cutoffs used in the proof.

The proof in Section~\ref{sec:log-rate} bounds the number $U_N$ of
unselected roots at $K=K_N$, in place of \eqref{eq:block-unselected}, by
\begin{equation}
 U_N\le N\Bigl[\frac1{K-1}+4\Bigl(\varepsilon_K^*(N)+\frac{D_K}N\Bigr)\Bigr]
 +Z_KN^{15/16}\log N+Z'_KN^{17/32}(\log N)^2,
 \label{eq:rate-unselected}
\end{equation}
outside the four radial events at the radii $1\pm(K_N-1)/N$ and
$1\pm K_N/N$ and the event of \eqref{eq:layer2-bad}.
Its probes are those of \eqref{eq:block-center-crossings}, at the fixed
width $2$, so that the probe error at the unit circle is at most
$\delta_2^*(N)$ and the crossing count at most $2\delta_2^*(N)N$.
By \eqref{eq:unit-block} at $x=0$, with $\delta_2^*(N)$ in place of
$\varepsilon_N$ as in Section~\ref{app:block-bounds}, and
\eqref{eq:rate-unselected} divided by $N$, the bound that we
use for the logarithmic rate is, for $N\le n\le N+m_j$,
\begin{align}
 \left|\frac{\nu_n(1)}n-\frac12\right|
 &\le3\delta_2^*(N)+\frac{U_N}N+\frac{2m_j}{N}\nonumber\\
 &\le3\Bigl(2N^{-1/64}+\frac2N
   +4\Bigl(\varepsilon_2^*(N)+\frac{D_2}N\Bigr)N^{1/64}\Bigr)
   +\frac1{K_N-1}\nonumber\\
 &\quad+4\Bigl(\varepsilon_{K_N}^*(N)+\frac{D_{K_N}}N\Bigr)
   +Z_{K_N}N^{-1/16}\log N\nonumber\\
 &\quad+Z'_{K_N}N^{-15/32}(\log N)^2+\frac{2m_j}{N}.
 \label{eq:E2-final}
\end{align}
Here $\varepsilon_{K_N}^*(N)$, $H_{K_N}=10^8e^{20(K_N+1)}$,
$D_{K_N}=10(K_N+1)^2e^{4K_N}$, $Z_{K_N}$ and $Z'_{K_N}$ are
$\varepsilon_K^*(N)$, $H_K$, $D_K$, $Z_K$ and $Z'_K$ at $K=K_N$, and
$D_2=90e^8$.
For this bound, the exceptional probability is \eqref{eq:block-failure}
with $K=K_N$, $R=2$ and $x=0$, where $C_{K_N}=V_{K_N}(D'_{K_N})^2$ and
$p_{K_N}^*(N)$ is the function $p_K^*(N)$ at $K=K_N$. By the third and
fourth lines of \eqref{eq:E2-majorants}, with $4H_0/256=H_0/64$, its terms
$4p_{K_N}^*(N)+C_{K_N}N^{-5/16}(\log N)^2$ are bounded as in the proof of
Theorem~\ref{thm:log-rate}.
By Lemma~\ref{lem:growing-width}, through \eqref{eq:E2-majorants}, every
term in \eqref{eq:E2-final} other than $1/(K_N-1)$ is $o(1/\log N)$, as
in Section~\ref{sec:log-rate}, since $K_N+1\le\log N$ and
$m_j/N\le9N^{-1/8}$.

\subsection{Nearby circles}\label{app:const-nearby}
Finally, we compute the constants of the proof of
Corollary~\ref{cor:crossings} in Section~\ref{sec:nearby-circles}.
For a target exponent $\kappa>0$, we choose a wider band with
\[
 0<\kappa_*<\min(\kappa,1/64),
\]
as does the choice $\kappa_*=\min(\kappa/2,1/128)$ in that proof.
Then the band $||z|-1|<N^{-1-\kappa_*}$ contains the circles
$|z|=1\pm N^{-1-\kappa}$, and $\kappa_*$ lies in the range
$0<\kappa\le1/64$ of the thin-band estimate of
Section~\ref{app:const-radial}.
We impose
\[
 N^{-1-\kappa}+2a/d_0<N^{-1-\kappa_*},
\]
where $2a/d_0$ is the radius of the selected disks of Step~2 of
Section~\ref{sec:layer3}. As in the proof of
Corollary~\ref{cor:crossings}, this condition places every good annular
root whose disk meets one of the circles in the band.
By the form of \eqref{eq:radial-count-near-unit} at $K=2$ in
Section~\ref{app:const-radial}, with $\kappa_*$ in place of $\kappa$,
\[
 \left|\frac{\nu_N(1+qN^{-1-\kappa_*})}N-\frac12\right|
 \le2N^{-\kappa_*}+4\varepsilon_2^*(N)N^{\kappa_*},\qquad q=-1,0,1,
\]
outside probability $5p_2^*(N)$.
Since $p_2^*(N)$ does not depend on the thin-band exponent, the
exceptional probability of the block is \eqref{eq:block-failure} with
$R=2$, where the term $5p_2^*(N)$ now bounds these five thin-band events
in place of the five probe events.
